\section{Analysis: H.264 as Adversarial Purifier}
\label{sec:analysis}

\subsection{The Codec-Kill Pattern}

Across all seven experiments, we observe a consistent pattern:
\begin{itemize}
  \item \textbf{Without codec}: attacks produce nonzero $\djfadv$ in at least some videos (mean $+0.005$ to $+0.121$), confirming that the adversarial objectives are meaningful and the optimizer finds solutions.
  \item \textbf{With CRF23}: $\djfauc$ collapses to near-zero in all cases (range: $[-0.004, +0.004]$). The codec reduces attack effectiveness by 96--100\%.
\end{itemize}

The attack strength without codec ($\djfadv$) has no predictive value for codec-robustness ($\djfauc$). The bus video in Mode F0-CD has the highest $\djfadv$ ($+0.867$) but $\djfauc = -0.002$. Videos with weaker $\djfadv$ are not more codec-robust.

\subsection{Hypothesized Root Cause: Frequency Concentration}

We propose the following mechanistic hypothesis to explain the codec-kill pattern. This remains a hypothesis---we do not provide direct spectral measurements---but it is grounded in well-known properties of both $L_\infty$ adversarial perturbations and DCT-based codecs.

Perturbations constrained to $\|\delta\|_\infty \leq \varepsilon$ in the pixel domain tend to concentrate energy in high spatial frequencies: to remain imperceptible against natural textures, they exploit neural network sensitivity at fine-grained spatial scales. H.264's DCT-based quantization at CRF23 zeroes out these high-frequency coefficients. Even though our feature-space attacks (Mode C+D) target a semantically deep representation (maskmem at $8 \times 8$ spatial resolution), the \emph{pixel-space} perturbation that drives these features is still subject to the $L_\infty$ budget and likely retains the same high-frequency structure.

This hypothesis predicts: (1) CRF18 would attenuate less (consistent with our CRF23 choice being a hard setting), (2) perturbations with explicit low-frequency structure (e.g., DCT-domain optimization) might survive better. We do not test these predictions experimentally, and the exact spectral decomposition of our perturbations is not measured. Future work should validate this mechanism directly via frequency-domain analysis.

\paragraph{Implication (if hypothesis holds).}
A codec-robust adversarial perturbation targeting SAM2 would require either: (1) increasing $\varepsilon$ sufficiently that low-frequency perturbation energy can meaningfully shift SAM2's memory state, or (2) a fundamentally different perturbation structure (e.g., DCT-domain or texture-level semantic manipulation) that does not rely on high-frequency patterns.

\subsection{Why Feature-Space Attacks Do Not Help}

One might expect that feature-space attacks targeting low-resolution representations (maskmem: $8 \times 8$) would produce pixel perturbations with more low-frequency content than pixel-space attacks. We observe this is not the case for $\varepsilon = 8/255$.

The gradient of the feature-space loss with respect to pixel $x_{ij}$ involves the Jacobian of the SAM2 encoder, which has large magnitude at high-spatial-frequency components (sharp edges, fine textures). Thus the optimal pixel perturbation $\delta_0^* = \arg\max J_\text{mem\_write}(\delta_0)$ subject to $\|\delta_0\|_\infty \leq \varepsilon$ tends to have the same high-frequency structure as direct pixel attacks, even though the \emph{objective} operates on low-resolution features.

\subsection{SAM2's Semantic Bottleneck}

A secondary factor may be SAM2's memory encoder architecture. The memory encoder compresses frame predictions to $8 \times 8$ tokens through a series of convolutions and attention operations. This spatial bottleneck naturally averages out high-frequency perturbation effects that survive codec at intermediate resolutions, providing a second layer of implicit purification beyond H.264 alone.

Together, H.264 CRF23 (primary) and SAM2's semantic bottleneck (secondary) make this threat model unusually resistant to pixel-constrained adversarial perturbations.
